\documentclass[11pt]{article}
\usepackage[utf8]{inputenc}
\usepackage[T1]{fontenc}
\usepackage[french]{babel}
\usepackage{graphicx}
\usepackage{geometry}
\usepackage{float}
\usepackage{listings}
\usepackage{xcolor}
\usepackage{booktabs}
\usepackage{array}
\usepackage{hyperref}
\usepackage{titlesec}
\usepackage{fancyhdr}
\usepackage{tikz}

\geometry{a4paper, margin=2.5cm}

% --- Palette de couleurs ---
\definecolor{eseoblue}{RGB}{0, 51, 102}
\definecolor{moderncyan}{RGB}{41, 128, 185}
\definecolor{darkgray}{RGB}{44, 62, 80}
\definecolor{lightgray}{RGB}{245, 247, 250}
\definecolor{codegray}{rgb}{0.5, 0.5, 0.5}
\definecolor{codepurple}{rgb}{0.58, 0, 0.82}
\definecolor{codeblue}{rgb}{0.0, 0.0, 0.8}
\definecolor{codegreen}{rgb}{0, 0.5, 0}
\definecolor{accentgreen}{RGB}{39, 174, 96}
\definecolor{accentred}{RGB}{192, 57, 43}

% --- Style des listings Java ---
\lstdefinestyle{javastyle}{
    backgroundcolor=\color{lightgray},
    commentstyle=\color{codegreen}\itshape,
    keywordstyle=\color{codeblue}\bfseries,
    numberstyle=\tiny\color{codegray},
    stringstyle=\color{codepurple},
    basicstyle=\ttfamily\footnotesize,
    breakatwhitespace=false,
    breaklines=true,
    captionpos=b,
    keepspaces=true,
    numbers=left,
    numbersep=5pt,
    showspaces=false,
    showstringspaces=false,
    showtabs=false,
    tabsize=2,
    language=Java,
    frame=single,
    rulecolor=\color{black!10}
}

\lstdefinestyle{plaintext}{
    backgroundcolor=\color{lightgray},
    basicstyle=\ttfamily\footnotesize,
    breaklines=true,
    frame=single,
    rulecolor=\color{black!10}
}

\lstset{style=javastyle}

% --- En-têtes ---
\pagestyle{fancy}
\fancyhf{}
\fancyhead[L]{\textit{STRMS -- Smart Task \& Resource Management System}}
\fancyhead[R]{\textit{E3e S6 -- Spring 2026}}
\fancyfoot[C]{\thepage}

\hypersetup{
    colorlinks=true,
    linkcolor=eseoblue,
    urlcolor=moderncyan,
    citecolor=eseoblue
}

\begin{document}

% ============================================================
% PAGE DE GARDE
% ============================================================
\begin{titlepage}
    \begin{tikzpicture}[remember picture, overlay]
        \fill[color=eseoblue] (current picture page.north west) rectangle ([xshift=1.2cm]current picture page.south west);
        \fill[color=moderncyan] ([xshift=1.2cm]current picture page.north west) rectangle ([xshift=1.35cm]current picture page.south west);
    \end{tikzpicture}

    \indent
    \begin{minipage}[t]{0.92\textwidth}
        \flushleft
        \vspace{-0.5cm}

        {\scshape\large\color{darkgray} ESEO -- Grande École d'Ingénieurs} \\
        {\small\color{gray} Département Informatique et Systèmes Embarqués} \\[0.2cm]
        {\small\color{gray} Année Universitaire 2025 -- 2026}

        \addvspace{4cm}

        {\fontsize{28}{34}\selectfont\bfseries\color{eseoblue} Smart Task \& Resource \\ Management System} \\[0.4cm]
        {\fontsize{16}{20}\selectfont\itshape\color{moderncyan} STRMS} \\[1.5cm]

        {\large\bfseries\color{darkgray} RAPPORT DE PROJET} \\
        {\normalsize\color{gray} Programmation Orientée Objet en Java}

        \addvspace{4cm}

        \begin{minipage}[t]{0.5\textwidth}
            \flushleft
            {\small\bfseries\color{eseoblue} RÉDIGÉ PAR :}\\[0.2cm]
            {\large\bfseries\color{darkgray} Romels \textsc{Ngongang}}\\
            {\large\bfseries\color{darkgray} Daniil \textsc{Vedishchev}}\\
            {\large\bfseries\color{darkgray} Kylan \textsc{Nom}}\\
            {\large\bfseries\color{darkgray} Nathan \textsc{Bensignor}}\\[0.4cm]
            {\small\color{gray} Promotion E3e -- Équipe N\textsuperscript{o} 4}
        \end{minipage}
        \hfill
        \begin{minipage}[t]{0.45\textwidth}
            \flushright
            {\small\bfseries\color{eseoblue} PRÉSENTÉ À :}\\[0.2cm]
            {\large\bfseries\color{darkgray} L'Équipe Pédagogique Java}\\[0.4cm]
            {\small\bfseries\color{eseoblue} DATE :}\\[0.1cm]
            {\large\color{darkgray} 4 juin 2026}
        \end{minipage}
    \end{minipage}
\end{titlepage}

\newpage
\setcounter{page}{1}
\tableofcontents
\newpage

% ============================================================
\section{Introduction}
% ============================================================

\subsection{Contexte du projet}
Ce rapport documente la conception et l'implémentation de l'application
\textbf{Smart Task \& Resource Management System (STRMS)}, projet réalisé dans
le cadre du module de \textit{Programmation Orientée Objet en Java} (E3e S6,
Spring 2026). Le sujet, défini par l'équipe pédagogique, consiste à concevoir
un système de gestion de tâches simulant un environnement d'ingénierie, avec
des contraintes fortes sur l'ordre d'exécution (dépendances), la sécurité
d'accès (rôles utilisateurs), la traçabilité (historique) et la persistance
(fichier). Le projet s'inscrit dans le programme du semestre S6, qui vise à
consolider les acquis de la programmation orientée objet à travers une
mise en situation réaliste et complète, depuis la conception UML jusqu'au
déploiement final.

Au-delà de l'application stricte des consignes, le projet a été l'occasion
de découvrir les outils modernes de l'écosystème Java : le système de
modules introduit par Java 9 (\textit{Java Platform Module System}), les
outils de packaging \texttt{jlink} et \texttt{jpackage}, ainsi que le
\textit{Maven Wrapper} pour reproduire l'environnement de compilation sur
n'importe quelle machine. Ces choix ont permis de livrer une application
non seulement fonctionnelle mais également distribuable, prête à être
testée par les membres du jury sans installation préalable de Maven ou de
JavaFX.

\subsection{Objectifs pédagogiques}
Le cahier des charges fixe dix objectifs pédagogiques auxquels le projet
répond intégralement. Le premier objectif est l'application des quatre
piliers de la programmation orientée objet : encapsulation, héritage,
polymorphisme et abstraction. Le deuxième concerne la conception et
l'interprétation de diagrammes UML, tant de classes que de séquence. Le
troisième objectif consiste à implémenter des classes abstraites et des
interfaces. Le quatrième porte sur l'utilisation systématique des
énumérations à la place de constantes brutes. Le cinquième implique la
manipulation des collections Java standard, notamment \texttt{ArrayList},
\texttt{HashMap}, \texttt{HashSet} et \texttt{PriorityQueue}. Le sixième
demande l'implémentation d'exceptions personnalisées avec une gestion
appropriée. Le septième porte sur la lecture et l'écriture de données dans
des fichiers. Le huitième est consacré à la gestion des dépendances entre
tâches et au maintien d'un historique d'audit. Le neuvième impose
l'écriture de tests unitaires et d'intégration avec JUnit. Enfin, le
dixième objectif porte sur l'organisation professionnelle d'un projet Java.

Chacun de ces objectifs sera revisité tout au long du présent rapport, en
montrant comment les choix d'implémentation y répondent concrètement. La
table de synthèse présentée en conclusion offre une vision panoramique de
cette correspondance.

\subsection{Composition de l'équipe et organisation}
Le projet a été réalisé en équipe de quatre étudiants. Romels Ngongang a
assuré la coordination, le modèle métier, la persistance fichier ainsi que
l'industrialisation du livrable (packaging \texttt{.exe} et \textit{Maven
Wrapper}). Daniil Vedishchev s'est concentré sur le contrôleur
\texttt{TaskManager} et la hiérarchie des exceptions personnalisées, qui
constituent le c\oe ur des règles métier. Kylan a développé l'interface
graphique JavaFX, depuis la conception de la charte graphique jusqu'à
l'intégration des différentes vues. Nathan Bensignor a rédigé les tests
JUnit, mis en place la couverture des scénarios obligatoires, et produit
la documentation technique distribuée avec le code. Le code source est
versionné sous Git, organisé selon une arborescence Maven standard, et
chaque commit reflète un incrément fonctionnel identifiable.

\subsection{Plan du document}
Le rapport est structuré pour suivre la logique de conception d'une
application orientée objet, de la vue d'ensemble vers les détails
d'implémentation. La section 2 décrit le système STRMS sous l'angle
fonctionnel. La section 3 expose l'architecture logicielle et les choix
d'organisation. Les sections 4, 5 et 6 documentent respectivement la
hiérarchie des utilisateurs, le modèle de tâche et le contrôleur central.
Les sections 7 et 8 traitent de la robustesse (exceptions) et de la
persistance. La section 9 présente l'interface graphique JavaFX, et la
section 10 la stratégie de validation par les tests JUnit. La section 11
décrit le packaging et la distribution multi-plateforme. Enfin, la section
12 dresse le bilan du projet et propose des perspectives d'évolution.

% ============================================================
\section{Présentation du système STRMS}
% ============================================================

\subsection{Fonctionnalités principales}
D'après le cahier des charges (section 3), le système doit permettre la
création d'utilisateurs avec des rôles spécifiques (Admin, Manager,
Engineer), ainsi que la création, l'assignation, la mise à jour et la
suppression de tâches. Il doit également suivre l'avancement des tâches et
leurs dépendances, maintenir un historique des tâches à des fins d'audit,
prioriser les tâches et gérer leur ordre d'exécution. Enfin, il doit
assurer le stockage des données du système dans des fichiers ainsi que la
génération de rapports et de tableaux de bord.

L'application implémente ces exigences à travers cinq vues principales
accessibles depuis une barre de navigation latérale. Le tableau de bord
synthétise des indicateurs clés tels que le nombre de tâches par statut, le
volume assigné à chaque ingénieur et la liste des tâches terminées. Le
tableau Kanban présente les tâches dans quatre colonnes correspondant aux
quatre états possibles, à la manière de l'outil industriel Jira. La vue
des utilisateurs récapitule la matrice de permissions par rôle. La vue
d'historique permet de consulter l'audit chronologique des actions
effectuées sur une tâche donnée. Enfin, la vue de rapports génère un
résumé textuel de l'état du système et propose des opérations de
sauvegarde et de chargement vers le format fichier dédié.

\subsection{Les trois rôles utilisateurs}
Le cahier des charges (section 3.3) définit trois rôles, chacun avec ses
propres responsabilités, et précise expressément que \textit{« seuls les
ingénieurs peuvent exécuter des tâches »}. Cette contrainte a guidé la
conception des permissions et a entraîné l'introduction d'une distinction
explicite entre l'exécution d'une tâche (au sens du démarrage) et sa
clôture (au sens de la validation finale).

\begin{table}[H]
\centering
\small
\begin{tabular}{p{2.5cm}p{12cm}}
\toprule
\textbf{Rôle} & \textbf{Responsabilités} \\
\midrule
\textbf{Admin}     & Tous les droits : créer, supprimer, assigner des tâches via le
                    \texttt{TaskManager}, et générer des rapports. Il représente
                    typiquement le chef de projet ou l'administrateur fonctionnel. \\
\textbf{Manager}   & Assigne les tâches aux ingénieurs, suit l'avancement et clôture
                    les tâches une fois validées. Génère des rapports. Il représente
                    le manager opérationnel d'une équipe d'ingénieurs. \\
\textbf{Engineer}  & Démarre les tâches qui lui sont assignées et met à jour leur
                    statut. Il est le seul rôle qui peut faire passer une tâche de
                    \texttt{TODO} à \texttt{IN\_PROGRESS}. En revanche, il ne peut
                    pas la clôturer lui-même : cette responsabilité est dévolue à un
                    Admin ou un Manager pour préserver un contrôle hiérarchique. \\
\bottomrule
\end{tabular}
\caption{Responsabilités des trois rôles utilisateurs.}
\end{table}

Notre interface permet de basculer en temps réel entre ces quatre
utilisateurs de démonstration, ce qui rend immédiatement visible l'effet
des permissions sur les actions disponibles. Les boutons inaccessibles
disparaissent ou apparaissent grisés selon le rôle courant, ce qui
constitue une démonstration vivante du polymorphisme appliqué au contrôle
d'accès.

\subsection{Le cycle de vie d'une tâche}
Une tâche évolue à travers quatre états définis par l'énumération
\texttt{TaskStatus}. Le cahier des charges (section 4.3.1) impose la
séquence de transitions \texttt{TODO} $\rightarrow$ \texttt{BLOCKED}
$\rightarrow$ \texttt{IN\_PROGRESS} $\rightarrow$ \texttt{DONE}.

\begin{figure}[H]
    \centering
    \begin{tikzpicture}[node distance=3cm, auto, thick,
        every node/.style={font=\small\sffamily},
        state/.style={rectangle, rounded corners=4pt, draw=eseoblue, fill=lightgray,
                      minimum width=2cm, minimum height=0.8cm, align=center}]
        \node[state] (todo)  {TODO};
        \node[state, right of=todo] (blocked)  {BLOCKED};
        \node[state, right of=blocked, xshift=0.5cm] (prog) {IN\_PROGRESS};
        \node[state, right of=prog, xshift=0.5cm] (done) {DONE};
        \draw[->] (todo) -- (blocked);
        \draw[->] (blocked) -- (prog);
        \draw[->] (prog) -- (done);
    \end{tikzpicture}
    \caption{Cycle de vie nominal d'une tâche.}
\end{figure}

L'état \texttt{TODO} correspond à une tâche fraîchement créée pour laquelle
le travail n'a pas commencé. Une tâche bascule en \texttt{BLOCKED} dès
qu'au moins une de ses dépendances n'est pas terminée, ce qui interdit son
exécution effective. L'état \texttt{IN\_PROGRESS} indique qu'un ingénieur
travaille actuellement dessus. Enfin, l'état \texttt{DONE} signale une
tâche achevée. Il s'agit d'un état \textbf{terminal} : aucune transition
n'en sort, ce qui garantit l'intégrité de l'historique d'exécution. Si une
correction est nécessaire après clôture, le cahier des charges précise
qu'une nouvelle tâche doit être créée, ce qui correspond à la pratique
industrielle d'éviter les régressions silencieuses.

Les transitions ne sont pas libres. La transition \texttt{TODO}
$\rightarrow$ \texttt{IN\_PROGRESS} ne peut s'effectuer que si toutes les
dépendances de la tâche sont au statut \texttt{DONE}. La transition
\texttt{IN\_PROGRESS} $\rightarrow$ \texttt{DONE} est, dans notre
implémentation, conditionnée à la permission \texttt{canCompleteTask()} de
l'acteur, ce qui implémente la règle hiérarchique évoquée plus haut.

\subsection{Scénario de démonstration}
Pour illustrer la dynamique du système, considérons le scénario suivant
inspiré directement du cahier des charges (section 3.3). Romels, qui
possède le rôle Admin, crée deux tâches : T1 (« Implémenter la page de
login ») et T2 (« Créer le schéma de base de données »). Il définit
ensuite une dépendance entre les deux : T2 dépend de T1. À cet instant,
T2 bascule automatiquement à l'état \texttt{BLOCKED} car sa dépendance
n'est pas encore terminée.

Daniil, qui est Manager, assigne T1 à Kylan, ingénieur. Lors de
l'assignation, T1 passe automatiquement à \texttt{IN\_PROGRESS} car elle
n'a pas de dépendance. Kylan travaille sur la tâche, puis la signale
comme prête à être validée. Daniil clôture T1, ce qui la fait passer à
\texttt{DONE}. Le système détecte alors que T2 n'a plus de dépendance non
satisfaite et la fait basculer automatiquement de \texttt{BLOCKED} à
\texttt{TODO}, où elle devient éligible à l'assignation. Toutes ces
opérations sont enregistrées dans l'historique de chaque tâche, ce qui
permet à n'importe quel utilisateur de tracer ultérieurement la chaîne
des décisions et des actions.

% ============================================================
\section{Architecture du code}
% ============================================================

\subsection{Structure en packages}
Le projet est organisé en six packages, chacun avec une responsabilité
précise. Cette séparation reflète les couches conceptuelles de
l'application : énumérations, exceptions, modèle, contrôleur, utilitaires
et interface graphique.

\begin{table}[H]
\centering
\small
\begin{tabular}{ll}
\toprule
\textbf{Package} & \textbf{Rôle} \\
\midrule
\texttt{fr.eseo.strms.enums}      & Les quatre énumérations du système \\
\texttt{fr.eseo.strms.exceptions} & Les sept exceptions personnalisées \\
\texttt{fr.eseo.strms.model}      & Modèle objet : User, Task, TaskHistoryEntry \\
\texttt{fr.eseo.strms.manager}    & Contrôleur central TaskManager \\
\texttt{fr.eseo.strms.utils}      & Persistance, rapports, notifications, tableau de bord \\
\texttt{fr.eseo.strms.ui}         & Interface graphique JavaFX \\
\bottomrule
\end{tabular}
\caption{Organisation des packages du projet.}
\end{table}

Cette ségrégation possède plusieurs vertus. Elle facilite la lecture du
code en permettant à un nouveau lecteur de localiser rapidement la classe
recherchée. Elle limite les dépendances cycliques, le modèle ne
connaissant pas les classes utilitaires et l'interface graphique étant
isolée du c\oe ur métier. Enfin, elle prépare une éventuelle évolution
du projet vers une architecture en modules au sens du \textit{Java Module
System}, ce qui a effectivement été réalisé pour les besoins du packaging.

\subsection{Modularité Java}
Le projet est déclaré comme un module Java explicite via le fichier
\texttt{module-info.java} placé à la racine de \texttt{src/main/java}.
Cette déclaration est indispensable pour pouvoir utiliser les outils
\texttt{jlink} et \texttt{jpackage} qui n'acceptent que des projets
modulaires. Elle déclare explicitement les dépendances du module (ici
\texttt{javafx.controls}, qui inclut transitivement \texttt{javafx.graphics}
et \texttt{javafx.base}), exporte le package contenant la classe
\texttt{Application} pour permettre à JavaFX de la charger, et ouvre les
packages dont les classes sont accédées par réflexion (en l'occurrence
\texttt{PropertyValueFactory} de \texttt{TableView}).

\subsection{Diagramme de classes}
Le diagramme statique met en évidence trois ensembles : la hiérarchie
d'utilisateurs avec ses trois sous-classes, le modèle de tâche avec
composition des dépendances et de l'historique, et le contrôleur central
qui orchestre les opérations.

\begin{figure}[H]
    \centering
    \includegraphics[width=\linewidth]{Diagrammedeclass.png}
    \caption{Diagramme de classes du système STRMS.}
\end{figure}

\subsection{Diagramme de séquence}
Le diagramme de séquence ci-dessous illustre l'ajout d'une dépendance et
la détection préventive de cycle, qui constitue l'interaction la plus
représentative du rôle central du \texttt{TaskManager}. Aucun composant
de l'interface graphique ne manipule directement le modèle : toute
modification passe par ce contrôleur, ce qui matérialise l'unicité de
l'autorité sur l'état du système.

\begin{figure}[H]
    \centering
    \includegraphics[width=0.85\linewidth]{Diagrammedeséquence.drawio.png}
    \caption{Diagramme de séquence : ajout d'une dépendance.}
\end{figure}

\subsection{Choix de conception transversaux}
Trois principes ont guidé l'ensemble des choix de conception. Le premier
est la centralisation des règles métier dans le \texttt{TaskManager}, qui
applique le patron de conception \textit{Facade} : aucune règle n'est
dispersée dans l'interface graphique ou dans le modèle. Le deuxième est
l'immuabilité partout où elle est pertinente, en particulier pour les
entrées d'historique qui doivent être inaltérables pour des raisons
d'audit. Le troisième est le polymorphisme à la place des tests sur le
type concret : le contrôleur ne contient aucun \texttt{instanceof} sur
les sous-classes de \texttt{User}, ce qui assure une extensibilité
naturelle (l'ajout d'un quatrième rôle ne nécessiterait aucune
modification du \texttt{TaskManager}).

% ============================================================
\section{La hiérarchie des utilisateurs}
% ============================================================

\subsection{La classe abstraite User}
Le cahier des charges (section 4.1) demande une classe \texttt{User}
abstraite avec des méthodes abstraites \texttt{canCreateTask()},
\texttt{canDeleteTask()}, etc. Cette classe matérialise simultanément
l'\textbf{abstraction} (on ne peut pas instancier un \texttt{User}
quelconque, il faut choisir un rôle), l'\textbf{encapsulation} (les
attributs sont privés et accédés uniquement via les getters et setters),
et l'\textbf{héritage} (les trois sous-classes étendent \texttt{User}).

\begin{lstlisting}[caption={Extrait de la classe abstraite User}]
public abstract class User {

    private final String id;
    private String name;
    private String email;

    protected User(String id, String name, String email) {
        this.id = id;
        this.name = name;
        this.email = email;
    }

    public abstract String getRole();
    public abstract boolean canCreateTask();
    public abstract boolean canDeleteTask();
    public abstract boolean canAssignTask();
    public abstract boolean canExecuteTask();
    public abstract boolean canCompleteTask();
    public abstract boolean canGenerateReport();
}
\end{lstlisting}

Le constructeur est déclaré \texttt{protected} pour interdire toute
instanciation directe de \texttt{User} : seules les sous-classes peuvent
en créer. Le champ \texttt{id} est \texttt{final}, ce qui signifie que
l'identifiant d'un utilisateur est fixé à la création et ne peut plus
changer pendant la durée de vie de l'objet, ce qui est conforme à la
sémantique d'un identifiant unique. Les six méthodes abstraites
définissent le contrat de permissions du système. Chaque sous-classe est
tenue de répondre à chacune de ces questions, ce qui constitue la base du
contrôle d'accès basé sur les rôles.

\subsection{Les sous-classes Admin, Manager et Engineer}
Chaque sous-classe implémente les six méthodes abstraites pour refléter
le rôle qu'elle représente. C'est un exemple direct de
\textbf{polymorphisme à liaison dynamique} : le code client appelle
uniformément \texttt{actor.canCreateTask()} sans connaître le type
concret de l'acteur, et chaque sous-classe répond pour elle-même au
moment de l'exécution.

\begin{lstlisting}[caption={Sous-classe Admin}]
public class Admin extends User {
    public Admin(String id, String name, String email) {
        super(id, name, email);
    }
    @Override public String  getRole()             { return "Admin"; }
    @Override public boolean canCreateTask()       { return true; }
    @Override public boolean canDeleteTask()       { return true; }
    @Override public boolean canAssignTask()       { return true; }
    @Override public boolean canExecuteTask()      { return false; }
    @Override public boolean canCompleteTask()     { return true; }
    @Override public boolean canGenerateReport()   { return true; }
}
\end{lstlisting}

Les classes \texttt{Manager} et \texttt{Engineer} suivent le même schéma
avec des combinaisons de \texttt{true} et \texttt{false} différentes
correspondant à leur rôle dans le système. La table suivante synthétise
cette matrice de permissions.

\begin{table}[H]
\centering
\small
\begin{tabular}{lccc}
\toprule
\textbf{Permission}            & \textbf{Admin} & \textbf{Manager} & \textbf{Engineer} \\
\midrule
\texttt{canCreateTask()}       & \textcolor{accentgreen}{\textbf{Oui}} & \textcolor{accentred}{\textbf{Non}} & \textcolor{accentred}{\textbf{Non}} \\
\texttt{canDeleteTask()}       & \textcolor{accentgreen}{\textbf{Oui}} & \textcolor{accentred}{\textbf{Non}} & \textcolor{accentred}{\textbf{Non}} \\
\texttt{canAssignTask()}       & \textcolor{accentgreen}{\textbf{Oui}} & \textcolor{accentgreen}{\textbf{Oui}} & \textcolor{accentred}{\textbf{Non}} \\
\texttt{canExecuteTask()}      & \textcolor{accentred}{\textbf{Non}}   & \textcolor{accentred}{\textbf{Non}} & \textcolor{accentgreen}{\textbf{Oui}} \\
\texttt{canCompleteTask()}     & \textcolor{accentgreen}{\textbf{Oui}} & \textcolor{accentgreen}{\textbf{Oui}} & \textcolor{accentred}{\textbf{Non}} \\
\texttt{canGenerateReport()}   & \textcolor{accentgreen}{\textbf{Oui}} & \textcolor{accentgreen}{\textbf{Oui}} & \textcolor{accentred}{\textbf{Non}} \\
\bottomrule
\end{tabular}
\caption{Matrice RBAC des permissions par rôle.}
\end{table}

\subsection{Implications de la matrice de permissions}
Cette matrice traduit une politique de séparation des rôles cohérente
avec la pratique industrielle. L'Admin concentre les droits de
configuration du système, ce qui correspond à un super-utilisateur ou à
un chef de projet. Le Manager partage avec l'Admin les droits
opérationnels d'assignation et de validation, mais ne peut pas modifier
la définition même des tâches (création et suppression). L'Engineer est
volontairement restreint à l'exécution opérationnelle, sans capacité de
modifier la structure du travail. La séparation entre
\texttt{canExecuteTask()} (réservé à l'Engineer) et
\texttt{canCompleteTask()} (réservé à l'Admin et au Manager) introduit
un point de contrôle hiérarchique : un ingénieur ne peut pas
auto-valider la fin d'une tâche, ce qui garantit une revue managériale
systématique. Lorsqu'un Engineer tente néanmoins de clôturer une tâche
depuis l'interface, le système affiche un message d'alerte explicite
expliquant la raison du refus et invitant à demander la validation au
manager.

% ============================================================
\section{Le modèle de tâche}
% ============================================================

\subsection{Les attributs de la classe Task}
Conformément au cahier des charges (section 4.2), la classe \texttt{Task}
encapsule ses propres données ainsi que les opérations qui agissent sur
ces données. Elle déclare neuf attributs privés couvrant l'identifiant,
le titre, la description, la priorité, le statut, la catégorie, la
deadline, l'ingénieur assigné, ainsi que deux listes : l'une pour les
dépendances vers d'autres tâches, l'autre pour l'historique des actions.

\begin{lstlisting}[caption={Attributs principaux de Task}]
public class Task implements Comparable<Task> {

    private final String id;
    private String title;
    private String description;
    private PriorityLevel priority;
    private TaskStatus status;
    private TaskCategory category;
    private LocalDate deadline;
    private Engineer assignedEngineer;

    private final List<Task> dependencies;
    private final List<TaskHistoryEntry> history;
    // ...
}
\end{lstlisting}

L'identifiant \texttt{id} est déclaré \texttt{final} pour les mêmes
raisons que dans la classe \texttt{User} : une tâche conserve son
identifiant à vie. Le statut initial est systématiquement \texttt{TODO},
ce qui correspond à l'état d'une tâche fraîchement créée. Les listes
\texttt{dependencies} et \texttt{history} sont déclarées \texttt{final},
ce qui signifie que la référence ne peut pas être réassignée, mais leur
contenu peut évoluer via les méthodes dédiées de la classe. À l'inverse,
elles ne sont jamais exposées telles quelles à l'extérieur : les
\texttt{getter} renvoient une vue non modifiable grâce à
\texttt{Collections.unmodifiableList}, ce qui empêche un code client
malveillant ou maladroit d'altérer la structure interne.

\subsection{La comparaison par priorité}
La classe \texttt{Task} implémente l'interface \texttt{Comparable<Task>},
ce qui autorise son utilisation directe dans une \texttt{PriorityQueue}
sans nécessiter de \texttt{Comparator} externe. La logique de comparaison
est centrale pour l'ordonnancement et mérite une attention particulière.

\begin{lstlisting}[caption={Comparaison par priorité}]
@Override
public int compareTo(Task other) {
    int cmp = Integer.compare(other.priority.getWeight(), this.priority.getWeight());
    if (cmp == 0) {
        return this.id.compareTo(other.id);
    }
    return cmp;
}
\end{lstlisting}

La méthode renvoie une valeur négative quand la tâche courante est plus
prioritaire que l'autre, ce qui place les tâches les plus critiques en
tête de la file. On notera l'inversion volontaire de l'ordre des
opérandes dans le \texttt{Integer.compare} : on compare le poids de
l'autre à celui de la tâche courante, et non l'inverse, ce qui réalise
le tri par priorité décroissante. En cas d'égalité de poids, on
départage par ordre lexicographique sur l'identifiant. Cette résolution
secondaire garantit la \textbf{stabilité} et la \textbf{reproductibilité}
de l'ordonnancement entre deux exécutions, ce qui est précieux à la fois
pour la prédictibilité du système et pour la rédaction de tests
déterministes.

\subsection{La classe TaskHistoryEntry et son immuabilité}
Le cahier des charges (section 6.2.3) exige que les entrées d'historique
soient \textbf{immuables} et en lecture seule après création. Cette
exigence n'est pas anodine : sans immuabilité, un acteur malveillant ou
un bug pourrait réécrire l'historique a posteriori, ce qui ruinerait la
valeur d'audit du système.

\begin{lstlisting}[caption={Classe TaskHistoryEntry, totalement immuable}]
public final class TaskHistoryEntry {

    private final String action;
    private final String performedBy;
    private final LocalDateTime timestamp;
    private final String description;

    public TaskHistoryEntry(String action, String performedBy, String description) {
        this(action, performedBy, description, LocalDateTime.now());
    }

    public TaskHistoryEntry(String action, String performedBy,
                            String description, LocalDateTime timestamp) {
        this.action = Objects.requireNonNull(action);
        this.performedBy = performedBy != null ? performedBy : "système";
        this.description = description != null ? description : "";
        this.timestamp = Objects.requireNonNull(timestamp);
    }

    public String getAction()           { return action; }
    public String getPerformedBy()      { return performedBy; }
    public LocalDateTime getTimestamp() { return timestamp; }
    public String getDescription()      { return description; }
}
\end{lstlisting}

Notre implémentation combine trois mécanismes complémentaires pour
garantir l'immuabilité. Premièrement, la classe est déclarée
\texttt{final}, ce qui interdit toute sous-classe susceptible de
modifier son comportement. Deuxièmement, l'ensemble des attributs sont
\texttt{final} et sont donc fixés au moment de la construction sans
possibilité de modification ultérieure. Troisièmement, la classe
n'expose aucun \texttt{setter}, ce qui élimine toute voie d'altération
externe. Une entrée d'historique, une fois inscrite, est ainsi conservée
telle quelle indéfiniment et peut servir de référence fiable pour un
audit ultérieur.

% ============================================================
\section{Le contrôleur central : TaskManager}
% ============================================================

\subsection{Présentation}
Le \texttt{TaskManager} est, selon le cahier des charges (section 3.2),
\textit{« le contrôleur central du système, responsable de toutes les
opérations sur les tâches et de l'application des règles métier »}. Il
agit comme un point d'entrée unique : aucune partie de l'application ne
manipule directement le modèle, tout passe par lui. Cette centralisation
garantit que les règles d'autorisation, l'intégrité du graphe de
dépendances et la traçabilité ne peuvent être contournées. Du point de
vue des patrons de conception, il s'agit d'une application du patron
\textit{Facade}.

\subsection{Les structures de données internes}
Les quatre attributs internes reflètent les recommandations du cahier
des charges (section 5).

\begin{lstlisting}[caption={Attributs internes du TaskManager}]
public class TaskManager {
    private final Map<String, Task> tasks;
    private final Map<String, User> users;
    private final Set<Task>  inProgressTasks;
    private final PriorityQueue<Task> readyQueue;

    public TaskManager() {
        this.tasks            = new HashMap<>();
        this.users            = new HashMap<>();
        this.inProgressTasks  = new HashSet<>();
        this.readyQueue       = new PriorityQueue<>();
    }
}
\end{lstlisting}

Le choix d'une \texttt{HashMap<String, Task>} permet une recherche d'une
tâche par identifiant en temps constant amorti, ce qui est essentiel
puisque l'opération de recherche est de loin la plus fréquente dans
l'application. La \texttt{HashMap<String, User>} joue un rôle analogue
pour les utilisateurs et permet une vérification instantanée des
permissions. Le \texttt{HashSet<Task>} suit les tâches actuellement en
cours d'exécution, sans doublon possible et avec un test d'appartenance
en temps constant. Enfin, la \texttt{PriorityQueue<Task>} maintient en
tête la tâche la plus prioritaire pouvant être exécutée, avec un coût
logarithmique pour les opérations d'insertion et de retrait. Cette
combinaison équilibre lisibilité, performance et respect des
recommandations pédagogiques.

\subsection{L'ajout d'une tâche}
La méthode \texttt{addTask} illustre le schéma général de toutes les
opérations du \texttt{TaskManager} : vérification des permissions de
l'acteur, validation des préconditions métier, insertion dans les
structures de données, puis enregistrement d'une trace dans l'historique.

\begin{lstlisting}[caption={Ajout d'une tâche}]
public void addTask(Task task, User actor)
        throws InvalidRoleException, DuplicateTaskException {

    ensurePermission(actor, actor.canCreateTask(), "créer une tâche");
    if (tasks.containsKey(task.getId())) {
        throw new DuplicateTaskException(
            "Une tâche avec l'identifiant '" + task.getId() + "' existe déjà.");
    }
    tasks.put(task.getId(), task);

    if (task.areDependenciesCompleted() && task.getStatus() == TaskStatus.TODO) {
        readyQueue.offer(task);
    }
    task.addHistoryEntry(new TaskHistoryEntry(
            "CREATION", actor.getName(),
            "Tâche '" + task.getTitle() + "' créée par " + actor.getRole()));
}
\end{lstlisting}

La méthode commence par s'assurer que l'utilisateur appelant possède le
droit de créer une tâche, sinon elle lève une
\texttt{InvalidRoleException}. Elle vérifie ensuite l'unicité de
l'identifiant, ce qui aurait déclenché une \texttt{DuplicateTaskException}
en cas de collision. La tâche est alors insérée dans la \texttt{HashMap}
principale. Si elle ne possède aucune dépendance, elle rejoint
immédiatement la file de priorité afin de pouvoir être ordonnancée.
Enfin, une entrée \texttt{CREATION} est enregistrée dans l'historique de
la tâche, ce qui assure la traçabilité de l'action.

\subsection{L'assignation d'une tâche}
L'assignation est réservée aux Admin et Manager. La méthode vérifie cette
permission, puis contrôle que la cible d'assignation est bien un
ingénieur. Si toutes les dépendances de la tâche sont satisfaites,
celle-ci passe directement à l'état \texttt{IN\_PROGRESS} et intègre
l'ensemble des tâches en cours, sinon elle reste à \texttt{BLOCKED}.

\begin{lstlisting}[caption={Assignation d'une tâche}]
public void assignTask(String taskId, String engineerId, User actor)
        throws TaskNotFoundException, InvalidRoleException, InvalidTaskStateException {

    ensurePermission(actor, actor.canAssignTask(), "assigner une tâche");
    Task task = requireTask(taskId);

    User user = users.get(engineerId);
    if (!(user instanceof Engineer engineer)) {
        throw new InvalidRoleException(
            "L'utilisateur " + engineerId + " n'est pas un ingénieur.");
    }
    task.setAssignedEngineer(engineer);
    task.addHistoryEntry(new TaskHistoryEntry(
            "ASSIGNMENT", actor.getName(),
            "Tâche assignée à " + engineer.getName()));

    if (task.areDependenciesCompleted() && task.getStatus() != TaskStatus.DONE) {
        task.updateStatus(TaskStatus.IN_PROGRESS);
        readyQueue.remove(task);
        inProgressTasks.add(task);
    } else if (task.getStatus() != TaskStatus.DONE) {
        task.updateStatus(TaskStatus.BLOCKED);
    }
}
\end{lstlisting}

On notera l'usage du \textit{pattern matching} sur \texttt{instanceof}
introduit par Java 16, qui combine en une seule ligne le test de type et
la conversion. Si l'utilisateur cible n'est pas un \texttt{Engineer}, une
\texttt{InvalidRoleException} est levée avec un message explicite.

\subsection{La clôture d'une tâche}
La méthode \texttt{completeTask} matérialise la règle métier introduite
selon laquelle seuls les Admin et Manager peuvent clôturer une tâche.
Elle effectue trois opérations essentielles. D'abord, elle vérifie que
toutes les dépendances sont satisfaites. Ensuite, elle marque la tâche
comme terminée et la retire des structures de suivi. Enfin, et c'est
l'étape la plus intéressante, elle propage l'événement : toutes les
tâches du système qui dépendaient de celle qui vient d'être clôturée
sont réévaluées.

\begin{lstlisting}[caption={Clôture d'une tâche avec propagation}]
public void completeTask(String taskId, User actor)
        throws TaskNotFoundException, InvalidRoleException,
               InvalidTaskStateException, DependencyNotCompletedException {

    ensurePermission(actor, actor.canCompleteTask(), "compléter une tâche");
    Task task = requireTask(taskId);

    if (!task.areDependenciesCompleted()) {
        throw new DependencyNotCompletedException(
            "Toutes les dépendances ne sont pas terminées.");
    }
    task.markAsDone();
    inProgressTasks.remove(task);
    readyQueue.remove(task);
    task.addHistoryEntry(new TaskHistoryEntry(
            "STATUS_CHANGE", actor.getName(),
            "Tâche terminée par " + actor.getRole()));

    // Propagation : toutes les taches qui dependaient de celle-ci
    for (Task other : tasks.values()) {
        if (other.getDependencies().contains(task)) {
            refreshTaskReadiness(other);
        }
    }
}
\end{lstlisting}

La méthode privée \texttt{refreshTaskReadiness} évalue chaque tâche
candidate à un déblocage. Si toutes ses dépendances sont désormais à
\texttt{DONE}, elle bascule de \texttt{BLOCKED} à \texttt{TODO} et
rejoint la file de priorité. Cette propagation en cascade est ce qui
permet au système de fonctionner sans intervention manuelle : l'utilisateur
n'a pas à se demander quelles tâches sont devenues éligibles, le système
les ramène automatiquement à la surface.

\subsection{La détection de cycles par DFS}
Le cahier des charges (section 6.1.2) impose explicitement un algorithme
de parcours de graphe pour détecter les cycles avant tout ajout de
dépendance. Le risque, sans détection, serait qu'une tâche se trouve
dans une situation d'interblocage où elle attend la complétion d'une
tâche qui elle-même l'attend.

Nous avons opté pour un parcours en profondeur (\textit{Depth-First
Search}) itératif utilisant une pile explicite. Le principe consiste à
se demander, avant d'ajouter l'arc « \texttt{task} dépend de
\texttt{dependsOn} », si \texttt{dependsOn} dépend déjà, directement ou
indirectement, de \texttt{task}. Si c'est le cas, l'ajout créerait un
cycle. On part donc de \texttt{dependsOn} et on suit ses dépendances en
remontant le graphe. Une pile explicite simule la récursivité, et un
\texttt{HashSet} mémorise les n\oe uds déjà visités pour éviter toute
boucle infinie sur un graphe contenant déjà des partages.

\begin{lstlisting}[caption={Détection de cycle par DFS itératif}]
public boolean detectCircularDependency(Task task, Task dependsOn) {
    Set<Task> visited = new HashSet<>();
    Deque<Task> stack = new ArrayDeque<>();
    stack.push(dependsOn);

    while (!stack.isEmpty()) {
        Task current = stack.pop();
        if (current.equals(task)) {
            return true;  // cycle detecte
        }
        if (visited.add(current)) {
            for (Task dep : current.getDependencies()) {
                stack.push(dep);
            }
        }
    }
    return false;
}
\end{lstlisting}

La complexité temporelle est en $O(V + E)$ dans le pire des cas, où $V$
représente le nombre de tâches et $E$ le nombre de dépendances : on
parcourt au plus une fois chaque tâche et chaque arc. La complexité
spatiale est en $O(V)$, induite par l'ensemble \texttt{visited} et la
pile. Le choix d'une implémentation itérative plutôt que récursive
évite tout risque de \texttt{StackOverflowError} sur des graphes très
profonds et offre un contrôle plus fin sur la consommation mémoire.

\subsection{La méthode addDependency et son comportement transactionnel}
L'ajout d'une dépendance fait appel à la détection de cycle. La méthode
\texttt{addDependency} effectue toutes les vérifications avant la moindre
mutation, de sorte qu'en cas de rejet, l'état du système reste
strictement identique. Cette propriété d'atomicité est essentielle pour
préserver l'intégrité du graphe en cas d'erreur, et elle correspond à
l'exigence du cahier des charges (section 6.3) selon laquelle
\textit{« les dépendances rejetées ne doivent pas altérer l'état
existant du système »}. Notons aussi qu'en cas de rejet, une trace
\texttt{DEPENDENCY\_REJECTED} est néanmoins ajoutée à l'historique de la
tâche cible, ce qui permet d'auditer ultérieurement les tentatives
infructueuses : la traçabilité prime sur l'absence de mutation
métier.

% ============================================================
\section{Les exceptions personnalisées}
% ============================================================

\subsection{Hiérarchie des exceptions}
Le cahier des charges (section 8) impose la définition de sept exceptions
personnalisées. Toutes héritent de \texttt{java.lang.Exception}, ce qui
les rend \textit{checked} : l'appelant est tenu de les traiter
explicitement soit par une clause \texttt{throws}, soit par un bloc
\texttt{try-catch}. Ce choix est délibéré et important pédagogiquement.
À l'opposé d'une \texttt{RuntimeException} qui peut être ignorée
silencieusement, une exception \textit{checked} oblige le compilateur à
vérifier la prise en charge, ce qui empêche tout silence sur un échec
métier.

\begin{table}[H]
\centering
\small
\begin{tabular}{lp{9.5cm}}
\toprule
\textbf{Exception} & \textbf{Condition de levée} \\
\midrule
\texttt{CircularDependencyException}     & Un ajout de dépendance créerait un cycle. \\
\texttt{DependencyNotCompletedException} & Démarrage ou clôture avec dépendances non terminées. \\
\texttt{TaskNotFoundException}           & Identifiant de tâche inconnu. \\
\texttt{InvalidRoleException}            & Action interdite par le rôle de l'utilisateur. \\
\texttt{FilePersistenceException}        & Erreur d'I/O lors de la sauvegarde ou du chargement. \\
\texttt{DuplicateTaskException}          & Création d'une tâche avec un identifiant existant. \\
\texttt{InvalidTaskStateException}       & Transition d'état invalide (ex. DONE vers IN\_PROGRESS). \\
\bottomrule
\end{tabular}
\caption{Les sept exceptions personnalisées du système.}
\end{table}

\subsection{Structure type d'une exception}
Toutes les exceptions personnalisées suivent une structure standard,
dont \texttt{CircularDependencyException} constitue un exemple
représentatif. Elles déclarent typiquement deux constructeurs : l'un
prenant uniquement un message, l'autre prenant message et cause pour
permettre l'enchaînement (\textit{exception chaining}) lorsque
l'exception métier enveloppe une exception technique sous-jacente.

\begin{lstlisting}[caption={Exemple : CircularDependencyException}]
public class CircularDependencyException extends Exception {
    public CircularDependencyException(String message) {
        super(message);
    }
    public CircularDependencyException(String message, Throwable cause) {
        super(message, cause);
    }
}
\end{lstlisting}

\subsection{Stratégie de remontée}
La stratégie adoptée consiste à laisser remonter les exceptions jusqu'à
la couche d'interface graphique, qui les capture et les transforme en
boîtes de dialogue d'erreur compréhensibles pour l'utilisateur final.
Cette approche garde la couche métier exempte de toute logique
d'affichage tout en assurant à l'utilisateur un retour clair en cas de
problème. Les exceptions \texttt{FilePersistenceException} en
particulier enveloppent les \texttt{IOException} sous-jacentes pour
ajouter un contexte métier, suivant le principe d'\textit{exception
translation}.

% ============================================================
\section{Persistance fichier}
% ============================================================

\subsection{Choix du format texte}
Plutôt que la sérialisation native de Java (\texttt{Serializable}), nous
avons opté pour un format texte personnalisé nommé \texttt{.strms}. Trois
raisons ont motivé ce choix. La première est la \textbf{lisibilité
humaine} : un fichier \texttt{.strms} peut être ouvert dans n'importe
quel éditeur pour vérification ou correction manuelle, ce qui est
précieux en phase de développement et de débogage. La deuxième est
l'\textbf{indépendance vis-à-vis de la JVM} : la sérialisation Java est
sensible aux changements de version et de classpath, alors qu'un fichier
texte est portable entre toutes les versions. La troisième est la
\textbf{robustesse à l'évolution} : si la structure d'une classe change,
le format binaire devient illisible, alors que le format texte peut être
migré ou tolérer des champs additionnels avec un peu de soin.

\subsection{Structure du format}
Chaque ligne du fichier est typée par un préfixe explicite. Le préfixe
\texttt{TASK} introduit la définition complète d'une tâche, suivi de ses
champs séparés par des points-virgules. Le préfixe \texttt{DEP} déclare
une dépendance entre deux tâches identifiées par leurs identifiants. Le
préfixe \texttt{HIST} enregistre une entrée d'historique. Le séparateur
\texttt{;} est échappé par \texttt{\textbackslash;} dans tout champ
contenant ce caractère, suivant un mécanisme classique d'échappement
avec \textit{placeholder} interne. Cette structure simple permet au
lecteur de comprendre instantanément un fichier de persistance.

\subsection{Algorithme de chargement avec traitement différé}
Lors d'une sauvegarde, l'ordre des lignes dans le fichier n'est pas
garanti car la méthode \texttt{HashMap.values()} ne préserve aucun ordre
particulier. Il est donc tout à fait possible qu'une ligne
\texttt{DEP;T-2;T-1} apparaisse dans le fichier avant la définition de
\texttt{TASK;T-1;...}. Un chargement naïf au fil de l'eau perdrait cette
dépendance puisque \texttt{T-1} ne serait pas encore présente dans la
\texttt{HashMap} au moment où l'on tenterait de la relier à \texttt{T-2}.

Notre solution consiste à effectuer une seule lecture du fichier en
chargeant immédiatement les lignes \texttt{TASK}, mais en
\textbf{accumulant les lignes \texttt{DEP} et \texttt{HIST} dans une
liste différée}. Une fois toutes les tâches chargées, on traite ces
lignes différées en parcourant la liste, et toutes les références sont
alors garanties d'être résolues correctement.

\begin{lstlisting}[caption={Chargement en une lecture avec traitement différé}]
Map<String, Task> loaded = new HashMap<>();
List<String[]> deferredLines = new ArrayList<>();

try (BufferedReader reader = new BufferedReader(new FileReader(filePath))) {
    String line;
    while ((line = reader.readLine()) != null) {
        if (line.isBlank()) continue;
        String[] parts = split(line);
        if ("TASK".equals(parts[0])) {
            loaded.put(parts[1], deserializeTask(parts, manager));
        } else {
            deferredLines.add(parts);
        }
    }
}

for (String[] parts : deferredLines) {
    switch (parts[0]) {
        case "DEP"  -> resolveDependency(loaded, parts);
        case "HIST" -> resolveHistoryEntry(loaded, parts);
    }
}
\end{lstlisting}

Cet algorithme est à la fois efficace (lecture séquentielle unique) et
robuste (immunisé contre les permutations de l'ordre des lignes). Il
illustre comment un problème de modélisation peut se traduire en
contraintes algorithmiques subtiles qui ne sont visibles qu'à
l'exécution sur des cas particuliers.

% ============================================================
\section{Interface graphique JavaFX}
% ============================================================

\subsection{Architecture des vues}
L'interface graphique est entièrement écrite en Java, sans recours à
FXML, ce qui permet un contrôle programmatique total et une lecture
linéaire du code d'interface. La fenêtre principale est organisée autour
d'une mise en page de type \texttt{BorderPane} composée de trois zones.
La zone supérieure contient une barre titre avec le nom de l'application
et un sélecteur d'utilisateur courant. La zone latérale gauche contient
une barre de navigation vers les cinq vues principales. La zone centrale
accueille la vue actuellement sélectionnée.

Les cinq vues principales sont les suivantes. La vue \texttt{Dashboard}
présente des cartes d'indicateurs colorées avec le nombre de tâches dans
chaque état, le détail des assignations par ingénieur et la liste des
tâches terminées. La vue \texttt{Kanban} affiche les tâches sous forme
de cartes dans quatre colonnes représentant les quatre états, à la
manière de l'outil industriel Jira. La vue \texttt{Users} synthétise les
permissions de chaque utilisateur dans un \texttt{TableView}. La vue
\texttt{History} affiche l'audit chronologique d'une tâche sélectionnée
via une \texttt{ComboBox}. La vue \texttt{Reports} génère un résumé
textuel et propose des actions de persistance fichier.

\subsection{Charte graphique}
La conception visuelle s'inspire d'applications industrielles modernes,
avec une dominante bleue (\texttt{\#0052CC}) et blanche, une barre
latérale sombre pour la navigation, des cartes blanches avec ombres
portées pour le contenu, et des badges colorés pour les niveaux de
priorité. L'ensemble est défini dans un fichier CSS unique
(\texttt{src/main/resources/css/style.css}), ce qui sépare proprement la
présentation de la logique.

\subsection{Adaptation dynamique aux permissions}
Le contenu de chaque vue s'adapte au rôle de l'utilisateur courant. Sur
le tableau Kanban, les boutons d'action sur chaque carte (Démarrer,
Terminer, Assigner, +Dépendance, Supprimer) apparaissent ou disparaissent
selon les permissions. Cette adaptation s'effectue uniquement à travers
les méthodes polymorphes \texttt{actor.canCreateTask()},
\texttt{actor.canCompleteTask()}, etc. Aucun test conditionnel sur le
type concret de l'acteur n'est jamais effectué, ce qui assure une
extensibilité naturelle.

Le cas particulier du bouton \textit{Terminer} pour un Engineer mérite
une mention. Plutôt que de simplement masquer le bouton, l'interface le
conserve visible mais grisé (style CSS \texttt{disabled}). Lorsque
l'Engineer clique dessus, une boîte de dialogue d'avertissement explique
que cette action est réservée aux Admin et Manager, et invite à demander
la validation au manager. Ce choix de conception privilégie la
\textbf{pédagogie} (l'utilisateur comprend pourquoi l'action est
indisponible) sur le simple masquage.

% ============================================================
\section{Tests automatisés JUnit}
% ============================================================

\subsection{Stratégie de test}
Le cahier des charges (section 7) impose une suite de tests isolés,
répétables et automatisés, validant à la fois les scénarios nominaux et
les cas d'erreur. Notre approche consiste à couvrir intégralement les
règles métier critiques du \texttt{TaskManager}, en s'appuyant sur deux
classes de test complémentaires.

\subsection{Structure}
La classe \texttt{TaskManagerTest} contient vingt-sept tests portant sur
les permissions par rôle, le cycle de vie complet des tâches, l'ajout et
le retrait de dépendances, la détection de cycles directs et indirects,
et le respect des transitions d'état. La classe
\texttt{FilePersistenceTest} contient deux tests d'intégration qui
valident le cycle complet de sauvegarde et de chargement, en vérifiant
que les dépendances et l'historique sont préservés à l'identique. Au
total, le projet contient \textbf{vingt-neuf tests}.

Chaque test est isolé via l'annotation \texttt{@BeforeEach} qui recrée
un \texttt{TaskManager} neuf, un Admin, un Manager et un Engineer. Cette
isolation garantit l'absence d'effet de bord entre tests, ce qui est
essentiel pour la fiabilité de la suite et permet l'exécution dans
n'importe quel ordre, conformément à l'exigence du cahier des charges
(section 7.3).

\subsection{Scénarios obligatoires couverts}
Le cahier des charges (section 7.2) liste cinq scénarios obligatoires.
Chacun est validé par un test dédié au nom explicite.

\begin{table}[H]
\centering
\small
\begin{tabular}{lp{6.5cm}}
\toprule
\textbf{Scénario obligatoire} & \textbf{Test JUnit correspondant} \\
\midrule
Ajout d'une dépendance valide        & \texttt{add\_valid\_dependency\_succeeds} \\
Rejet d'une dépendance circulaire    & \texttt{direct\_circular\_dependency\_is\_rejected} \\
Rejet circulaire indirect            & \texttt{indirect\_circular\_dependency\_is\_rejected} \\
Intégrité du graphe après rejet     & \texttt{graph\_integrity\_preserved\_after\_rejection} \\
Suppression correcte de dépendance  & \texttt{remove\_dependency\_works} \\
\bottomrule
\end{tabular}
\caption{Couverture des scénarios obligatoires par les tests JUnit.}
\end{table}

\subsection{Exemple représentatif}
Le test le plus représentatif est celui qui valide le rejet d'un cycle
indirect. Il prépare un graphe acyclique \texttt{C $\to$ B $\to$ A},
puis tente d'ajouter l'arc \texttt{A $\to$ C} qui créerait un cycle
indirect.

\begin{lstlisting}[caption={Test du rejet d'une dépendance circulaire indirecte}]
@Test
@DisplayName("Rejet d'une dépendance circulaire indirecte")
void indirect_circular_dependency_is_rejected() throws Exception {
    manager.addTask(makeTask("A", PriorityLevel.HIGH), admin);
    manager.addTask(makeTask("B", PriorityLevel.HIGH), admin);
    manager.addTask(makeTask("C", PriorityLevel.HIGH), admin);

    manager.addDependency("B", "A", admin); // B -> A
    manager.addDependency("C", "B", admin); // C -> B

    // Tentative : A -> C  produirait  A -> C -> B -> A (cycle)
    assertThrows(CircularDependencyException.class,
            () -> manager.addDependency("A", "C", admin));
}
\end{lstlisting}

La méthode \texttt{assertThrows} de JUnit vérifie qu'une
\texttt{CircularDependencyException} est bien levée par l'opération
d'ajout, conformément au comportement attendu. L'utilisation de
l'annotation \texttt{@DisplayName} fournit un libellé lisible dans les
rapports d'exécution, ce qui facilite l'analyse en cas d'échec.

\subsection{Résultats et couverture}
L'ensemble des vingt-neuf tests s'exécute en moins d'une seconde et
passe sans aucune erreur sur les machines des quatre membres de
l'équipe, ce qui valide à la fois la robustesse du code et la
reproductibilité de l'environnement de build via le \textit{Maven Wrapper}.
Les tests couvrent toutes les exceptions personnalisées : chacune est
levée et capturée au moins une fois dans la suite, ce qui garantit que
les chemins d'erreur ne sont pas du code mort.

% ============================================================
\section{Packaging et distribution}
% ============================================================

\subsection{Enjeu}
Au-delà de la mise au point du code, nous avons souhaité produire une
distribution facile à tester pour le jury et nos camarades. L'expérience
courante est qu'une application Java « simple à lancer » se transforme
souvent en cauchemar quand l'utilisateur ne dispose pas du bon JDK ou
n'a pas configuré JavaFX. Pour anticiper ce problème, nous avons mis en
place deux modes de distribution complémentaires.

\subsection{Mode universel : Maven Wrapper}
Le projet embarque le \textit{Maven Wrapper} (\texttt{mvnw},
\texttt{mvnw.cmd}, dossier \texttt{.mvn/wrapper}). Ces scripts
téléchargent automatiquement la version exacte de Maven attendue par le
projet (3.9.9) au premier lancement, ce qui évite à l'utilisateur
d'installer Maven manuellement. Le seul prérequis devient un JDK 17 ou
supérieur. Maven se charge ensuite du téléchargement des dépendances
JavaFX et de la sélection de la variante adaptée à l'architecture (x64,
ARM Apple Silicon, etc.). La commande à exécuter est minimale :
\texttt{mvnw javafx:run} sur Windows ou \texttt{./mvnw javafx:run} sur
macOS et Linux.

\subsection{Mode autonome : installeur Windows via jpackage}
Pour les utilisateurs ne souhaitant rien installer du tout, nous avons
produit un installeur Windows à l'aide des outils \texttt{jlink} et
\texttt{jpackage}, tous deux inclus dans le JDK 17 depuis Java 14.
L'outil \texttt{jlink} produit une image runtime contenant uniquement
les modules Java strictement nécessaires à l'application :
\texttt{java.base}, \texttt{java.desktop}, \texttt{javafx.controls},
\textit{etc}. L'outil \texttt{jpackage} encapsule ensuite cette image,
les classes de l'application et un launcher \texttt{STRMS.exe} dans un
dossier autonome de 91 Mo environ.

Sur Windows 11, la fonctionnalité \textit{Smart App Control} peut bloquer
les binaires \texttt{.exe} non signés numériquement. Pour contourner
cette limitation sans recourir à un certificat de signature payant, le
script de packaging produit en parallèle un launcher \texttt{STRMS.bat}
qui invoque directement l'exécutable Java embarqué dans le dossier
\texttt{runtime}. Le \texttt{.bat} n'étant pas filtré par Smart App
Control, il fonctionne sans avertissement sur toutes les machines
Windows.

\subsection{Bénéfices observés}
La coexistence des deux modes a permis de tester l'application sur les
machines de tous les membres de l'équipe, y compris celles équipées de
Mac M1, sans aucune intervention manuelle particulière. Le mode autonome
a également permis de partager l'application avec des étudiants
extérieurs à l'équipe pour des retours rapides, ce qui aurait été
nettement plus difficile avec une distribution classique sous forme de
code source à compiler.

% ============================================================
\section{Conclusion}
% ============================================================

\subsection{Bilan des objectifs pédagogiques}
Le projet STRMS valide l'ensemble des objectifs pédagogiques fixés par le
cahier des charges. Le tableau ci-dessous synthétise la correspondance
entre chaque concept demandé et sa mise en \oe uvre concrète.

\begin{table}[H]
\centering
\small
\begin{tabular}{lp{8.5cm}}
\toprule
\textbf{Concept} & \textbf{Mise en \oe uvre concrète dans STRMS} \\
\midrule
Encapsulation     & Attributs privés, listes exposées en lecture seule via \texttt{Collections.unmodifiableList}. \\
Héritage          & \texttt{User} abstract avec trois sous-classes concrètes. \\
Polymorphisme     & Appels à \texttt{actor.canCreateTask()} sans aucun \texttt{instanceof}. \\
Abstraction       & Classe \texttt{User} \texttt{abstract} et interface \texttt{Reportable}. \\
Énumérations      & Quatre enums utilisées partout au lieu de constantes brutes. \\
Collections       & \texttt{HashMap}, \texttt{HashSet}, \texttt{PriorityQueue}, \texttt{ArrayList}. \\
Exceptions        & Sept exceptions personnalisées, toutes \textit{checked}. \\
File I/O          & Format \texttt{.strms} avec sauvegarde et chargement. \\
JUnit             & Vingt-neuf tests, scénarios obligatoires couverts. \\
UML               & Diagramme de classes et diagramme de séquence. \\
\bottomrule
\end{tabular}
\caption{Bilan des concepts mis en \oe uvre.}
\end{table}

Au-delà des concepts demandés explicitement, le projet a permis d'explorer
des thèmes transversaux importants pour la pratique professionnelle de
l'ingénierie logicielle : la modélisation d'un graphe orienté acyclique
appliquée à un problème métier concret, la conception d'une interface
graphique réactive et lisible, la définition d'un format de persistance
texte robuste, et la stratégie d'industrialisation d'un livrable.

\subsection{Difficultés rencontrées et solutions}
Trois difficultés méritent d'être mentionnées. La première est l'ordre
non garanti lors de la sérialisation. Notre première version du
chargement perdait silencieusement certaines dépendances car les lignes
\texttt{DEP} pouvaient apparaître avant les \texttt{TASK} référencées
dans le fichier. La solution adoptée consiste en un chargement en une
seule lecture avec traitement différé des lignes \texttt{DEP} et
\texttt{HIST}, comme expliqué dans la section consacrée à la persistance.
Cette solution est efficace, élégante et préserve les performances.

La deuxième difficulté concernait la définition stricte des transitions
d'état avec \texttt{DONE} comme état terminal. Identifier précisément
les transitions autorisées et celles à interdire a nécessité un travail
de modélisation important. La solution a été centralisée dans la
méthode \texttt{updateStatus} de la classe \texttt{Task}, qui rejette
toute transition sortant de \texttt{DONE} par une
\texttt{InvalidTaskStateException} explicite.

La troisième difficulté est apparue lors du packaging Windows. Le
\texttt{STRMS.exe} produit par \texttt{jpackage} étant non signé
numériquement, il était bloqué par Smart App Control de Windows 11. La
solution adoptée a été de produire en parallèle un launcher
\texttt{STRMS.bat} qui invoque directement l'exécutable Java embarqué,
contournant la vérification de signature.

\subsection{Perspectives}
Plusieurs pistes ont été identifiées pour faire évoluer le projet vers
une solution de qualité industrielle. La première serait la migration
du format \texttt{.strms} vers une base de données relationnelle, par
exemple SQLite ou PostgreSQL via JDBC, ce qui apporterait une meilleure
gestion de la concurrence et une montée en charge plus aisée. La
deuxième consisterait à ajouter un véritable système d'authentification
login et mot de passe avec hachage par BCrypt, à la place du sélecteur
d'utilisateur actuel. La troisième serait l'implémentation effective des
canaux de notification \texttt{EMAIL} et \texttt{SMS}, aujourd'hui
simulés sur la console. La quatrième piste serait de compléter les
tests JUnit par des tests d'interface graphique automatisés à l'aide
d'un framework comme TestFX. Enfin, une cinquième évolution intéressante
serait l'exposition des fonctionnalités du \texttt{TaskManager} sous
forme d'une API REST, ce qui permettrait à plusieurs clients de
travailler simultanément sur le même ensemble de tâches.

\subsection{Conclusion finale}
Le projet STRMS, en plus de couvrir intégralement les exigences
pédagogiques du cahier des charges, illustre la mise en pratique
conjointe des piliers de la programmation orientée objet et d'un
algorithme classique de la théorie des graphes appliqué à un cas
industriel concret. L'application est entièrement fonctionnelle, validée
par vingt-neuf tests automatisés, distribuée sous deux formats
complémentaires et accompagnée d'une documentation complète facilitant
sa reprise par toute personne extérieure au projet. Le travail mené à
cette occasion a constitué une mise en situation enrichissante,
fidèle à ce qu'un ingénieur logiciel rencontrera dans sa pratique
professionnelle quotidienne.

\end{document}
